%% CV - Oda, Senior Data Analyst (vikariat)
%% Compile with: lualatex -interaction=nonstopmode main_oda.tex

\documentclass[11pt,a4paper,sans]{moderncv}
\moderncvstyle{banking}
\moderncvcolor{blue}

\renewcommand*{\firstnamestyle}[1]{{\fontsize{34}{36}\bfseries\upshape\color{color1}#1}}
\renewcommand*{\lastnamestyle}[1]{{\fontsize{34}{36}\bfseries\upshape\color{color1}#1}}
\renewcommand*{\sectionstyle}[1]{{\sectionfont\color{color1}#1}}

\usepackage[utf8]{inputenc}
\usepackage{needspace}
\usepackage{hyperref}
\hypersetup{
    colorlinks=true,
    linkcolor=blue,
    urlcolor=blue,
    pdftitle={Ashbin Jaison - CV},
    pdfpagemode=FullScreen,
}
\usepackage[scale=0.79]{geometry}
\nopagenumbers{}

\name{Ashbin}{Jaison}
\address{Bergen, Norway (open to relocation)}{}{}
\phone[mobile]{+47 486 251 21}
\email{ashbinjaison@gmail.com}
\extrainfo{\href{https://www.linkedin.com/in/ashbin-jaison/}{LinkedIn}, \href{https://github.com/ashbin-jaison}{GitHub}}

\begin{document}

\makecvtitle
\enlargethispage{8\baselineskip}

% ============================================================
%     PROFILE STATEMENT
% ============================================================

\vspace{2pt}
\small{Quantitative data scientist with a PhD in applied statistics and 5 years building analytical frameworks, defining impact metrics, and translating complex datasets into decisions for cross-functional stakeholders. Core stack is Python and SQL; experienced in statistical evaluation, hypothesis testing, A/B testing methodology, and designing measurement frameworks in fast-moving, iterative product environments.}

% ============================================================
%     CORE COMPETENCIES
% ============================================================

\section{Core Competencies}
\vspace{1pt}
\begin{itemize}

\item \textbf{Data Analytics \& Python}: pandas, NumPy, scikit-learn, XGBoost, LightGBM, matplotlib; large dataset analysis; statistical modelling and hypothesis testing to measure outcome improvements

\item \textbf{Impact Metrics \& Experimentation}: Defined and tracked success metrics across a 4-year programme; designed controlled comparison experiments and interpreted statistical results to measure outcome improvements; probabilistic evaluation (RMSE, Brier score, reliability analysis); systematic model calibration and refinement

\item \textbf{SQL \& Data Modelling}: SQL for analytical queries, data extraction, and data modelling; automated data ingestion and processing pipelines (ECMWF API, NetCDF/GRIB, xarray) for operational datasets

\item \textbf{Stakeholder Communication \& Visualisation}: Translating complex quantitative outputs into accessible insights, reports, and dashboards (Python/matplotlib, Power BI); bridging data analysis and business decisions for technical and non-technical audiences; improving analytical approaches for non-technical stakeholders

\end{itemize}

% ============================================================
%     PROFESSIONAL EXPERIENCE
% ============================================================

\section{Professional Experience}
\vspace{3pt}
\begin{itemize}

\needspace{5\baselineskip}
\item{\cventry{Mar 2025--Jan 2026}{Climate Data Scientist}{Admaren Tech.}{Cochin, India}{}{\vspace{1pt}
\begin{itemize}
    \item Built automated data pipelines and delivered analytics on operational ocean-weather datasets; collaborated with software engineers to integrate outputs into production systems
    \item Quantified extreme event probabilities to support maritime risk decisions; produced dashboards and analytical reports translating findings into operational recommendations for clients
\end{itemize}}}

\vspace{3pt}

\needspace{5\baselineskip}
\item{\cventry{Jan 2021--Jan 2025}{PhD Research Fellow, Climate Risk}{University of Bergen}{Bergen, Norway}{}{\vspace{1pt}
\begin{itemize}
    \item Defined impact and success metrics for windstorm prediction models; designed evaluation frameworks and interpreted experimental results to measure outcome improvements; improved forecast quality by $\sim$40\% through systematic hypothesis testing and calibration against observed outcomes
    \item Addressed ambiguous analysis problems by gathering domain expert input and synthesising coherent evaluation frameworks; identified and corrected flawed model assumptions to improve output reliability
    \item Produced analytical reports and risk insights for insurance stakeholders; translated complex probabilistic outputs into accessible findings and recommendations, enabling non-technical decision-makers to act on model results
\end{itemize}}}

\end{itemize}

% ============================================================
%     INDEPENDENT PROJECTS
% ============================================================

\section{Selected Projects}
\vspace{1pt}
\begin{itemize}
\item \textbf{Wind Power Forecasting -- NORA AI Hackathon (2025):} Built and benchmarked ensemble ML models (XGBoost, LightGBM, Random Forest); improved forecast accuracy by $>$50\% via systematic feature engineering and optimisation. Ran structured model comparisons as controlled experiments with defined metrics; documented methodology, results, and evaluation framework.
\end{itemize}

% ============================================================
%     EDUCATION
% ============================================================

\section{Education}
\vspace{1pt}
\begin{itemize}

\item{\cventry{Jan 2021--Jan 2025}{PhD, Applied Statistics -- Climate Risk Modeling}{University of Bergen}{Bergen, Norway}{}{}}

\vspace{3pt}

\item{\cventry{2018--2020}{M.Sc Statistics (First Class with Distinction)}{Cochin University of Science and Technology}{India}{}{}}

\end{itemize}

% ============================================================
%     LANGUAGES & CERTIFICATIONS
% ============================================================

\section{Languages \& Certifications}
\vspace{1pt}
\begin{itemize}
\item English (fluent), Norwegian (beginner)
\item Introduction to Generative AI -- Google Cloud (2025); Azure Fundamentals -- DataCamp (2025); Machine Learning Scientist -- DataCamp; Deep Learning Professional Certificate -- IBM/edX
\end{itemize}

\end{document}
